
\documentclass[12pt]{article}

\usepackage{geometry}
\usepackage{titling}
\usepackage{amsmath,amssymb}
\usepackage{fontspec}
\usepackage{unicode-math}
\usepackage{array}
\usepackage{colortbl}
\usepackage[dvipsnames]{xcolor}
\usepackage[most]{tcolorbox}

\geometry{a4paper, top=2.5cm, bottom=2.5cm, left=2.6cm, right=2.5cm}
\setlength{\droptitle}{-3cm}

\setmonofont{DejaVu Sans Mono}
\setmainfont{Libertinus Serif}
\setsansfont{Libertinus Sans}
\setmathfont{STIXTwoMath-Regular.otf}

\newcommand{\MVec}[1]{\mathbf{#1}}

\title{Application of Euclidean Distance \\ in \\ K-Nearest Neighbor Regression}
\author{Devi Prasad M}
\date{August 2026}

\begin{document}
\pagecolor{yellow!30!brown!40}

\maketitle

\section*{Euclidean Distance -- An Application: KNN Regression}

\subsection*{A 4D Example: Predicting House Prices}
Consider a 4-dimensional feature space where we want to predict
the sale price of a house based on its characteristics.

Our features are:
\begin{itemize}
    \item $x^{(1)}$: Living Area (in $1,000$ sq ft) -- \textit{Continuous}
    \item $x^{(2)}$: Distance to City Center (kms) -- \textit{Continuous}
    \item $x^{(3)}$: Age of Property (years) -- \textit{Integer}
    \item $x^{(4)}$: Number of Bedrooms -- \textit{Integer}
\end{itemize}

Our target $y$ is the \textbf{Sale Price (in INR Crores)}.

\vspace{0.5cm}
\noindent Suppose a new house comes on the market. Our query point
is a $2,000$ sq ft home, $1.0$ km from the center, $10$ years old,
with $3$ bedrooms:
\[
    \MVec{x}_q = (2.0, 1.0, 10, 3)
\]

\noindent Suppose our historical training dataset contains the
following $12$ recent sales:
\[
    \begin{array}{c|cccc|c}
        \text{Property} &
            x^{(1)} \text{ (Area)} &
            x^{(2)} \text{ (Distance)} &
            x^{(3)} \text{ (Age)} &
            x^{(4)} \text{ (Beds)} &
            y_i \text{ (Price)} \\ \hline
        \MVec{x}_1 & 1.5 & 1.5 & 12 & 3 & 3.5 \\
        \MVec{x}_2 & 2.2 & 0.5 & 8  & 4 & 5.2 \\
        \MVec{x}_3 & 2.0 & 1.0 & 9  & 3 & 4.8 \\
        \MVec{x}_4 & 1.8 & 1.2 & 10 & 3 & 4.6 \\
        \MVec{x}_5 & 2.5 & 2.0 & 5  & 4 & 6.0 \\
        \MVec{x}_6 & 1.2 & 3.0 & 20 & 2 & 2.5 \\
        \MVec{x}_7 & 2.1 & 0.8 & 11 & 3 & 4.9 \\
        \MVec{x}_8 & 3.0 & 5.0 & 2  & 5 & 7.0 \\
        \MVec{x}_9 & 1.9 & 1.5 & 10 & 2 & 4.1 \\
        \MVec{x}_{10} & 2.0 & 0.5 & 15 & 3 & 4.3 \\
        \MVec{x}_{11} & 2.3 & 1.1 & 10 & 4 & 5.4 \\
        \MVec{x}_{12} & 1.7 & 1.0 & 8  & 3 & 4.4
    \end{array}
\]
We choose $K=3$.

\subsection*{Step 1: Measure the Euclidean Distance to Every Point}
The 4-dimensional Euclidean distance to any training
point $\MVec{x}_i$ is:
\[
    \mathbf{dist}(\MVec{x}_q,\MVec{x}_i)
    = \sqrt{
        (2.0-x_i^{(1)})^2 +
        (1.0-x_i^{(2)})^2 +
        (10-x_i^{(3)})^2  +
        (3-x_i^{(4)})^2
    }
\]

We calculate this Euclidean distance to \textit{every} one of the $12$ properties:
\[
    \begin{array}{c|c|c@{\qquad}c|c|c}
        \text{Property} & y_i \text{ (Price)} & \mathbf{dist}(\MVec{x}_q,\MVec{x}_i)
        & \text{Property} & y_i \text{ (Price)} & \mathbf{dist}(\MVec{x}_q,\MVec{x}_i) \\ \hline
        \MVec{x}_1 & 3.5 & 2.12 & \MVec{x}_7 & 4.9 & 1.02 \\
        \MVec{x}_2 & 5.2 & 2.30 & \MVec{x}_8 & 7.0 & 9.22 \\
        \MVec{x}_3 & 4.8 & 1.00 & \MVec{x}_9 & 4.1 & 1.12 \\
        \MVec{x}_4 & 4.6 & 0.2828 & \MVec{x}_{10} & 4.3 & 5.02 \\
        \MVec{x}_5 & 6.0 & 5.22 & \MVec{x}_{11} & 5.4 & 1.05 \\
        \MVec{x}_6 & 2.5 & 10.28 & \MVec{x}_{12} & 4.4 & 2.02
    \end{array}
\]

\[
    \boxed{
    \text{Measure: Every training point is compared with the query.}
    }
\]

\subsection*{Step 2: Select the $K$ Nearest Neighbors}
Since $K=3$, we select the three properties with the smallest
Euclidean distances:
\[
    \begin{array}{c|c|c}
        \text{Neighbor} & \text{Target } y_i & \mathbf{dist}(\MVec{x}_q,\MVec{x}_i) \\ \hline
        \MVec{x}_4=(1.8, 1.2, 10, 3) & 4.6 & 0.2828 \\
        \MVec{x}_3=(2.0, 1.0, 9, 3) & 4.8 & 1.00 \\
        \MVec{x}_7=(2.1, 0.8, 11, 3) & 4.9 & 1.02
    \end{array}
\]

Therefore, our nearest neighborhood is:
\[
    N_3(\MVec{x}_q) =
    \left\{
        (\MVec{x}_4, 4.6),
        (\MVec{x}_3, 4.8),
        (\MVec{x}_7, 4.9)
    \right\}.
\]

\subsection*{Step 3: Predict}
The KNN regression prediction is the arithmetic mean of the target
prices of these three most similar houses:
\[
    \hat{y}(\MVec{x}_q) = \frac{4.6 + 4.8 + 4.9}{3} = \frac{14.3}{3} \approx 4.77
\]

The KNN estimate for the unknown market value of our
query house is:
\[
    \boxed{
        \text{Estimated Price } \approx \text{INR } 4.77 \text{ Crores}
    }
\]

\newpage

\subsection*{A Critical Flaw: The Need for Feature Scaling}
While the example above perfectly demonstrates the mechanics of KNN,
there is a hidden mathematical flaw in applying Euclidean distance
directly to raw, unscaled real-world data.

Consider the differences in scale among our features:
\begin{itemize}
    \item A difference of $5$ years in \textbf{Age}
          contributes $(5)^2 = 25$ to the distance sum.
    \item A difference of $2$ kms in \textbf{Distance}
          contributes only $(2)^2 = 4$ to the distance sum.
\end{itemize}

Even if distance to the city center has a more massive impact on
actual property value, the \textbf{Age} feature mathematically
dominates the Euclidean distance simply because its numerical
values are larger. In effect, we are comparing apples to oranges.

\vspace{0.3cm}
\noindent\textbf{The Solution: Standardization} \\
To prevent variables with larger scales from biasing the distance
metric, distance-based algorithms require data preprocessing.
We must normalize or standardize the features so they
contribute equally to the geometry.

A standard approach is \textit{standardization} or
\textit{Z-score normalization},
which transforms every feature to have a mean of $0$ and a
standard deviation of $1$:
\[
    x^{(j)}_{\text{scaled}} = \frac{x^{(j)} - \mu_j}{\sigma_j}
\]
where $\mu_j$ is the mean and $\sigma_j$ is the standard deviation of feature $j$.


\newpage

\section*{Z-Score Normalization and Re-computation}

\subsection*{Step 4: Normalizing the Dataset (Z-score Normalization)}
\[
    \mu_j = \frac{1}{N} \sum_{i=1}^{N} x_i^{(j)},
    \qquad
    \sigma_j = \sqrt{\frac{1}{N} \sum_{i=1}^{N} (x_i^{(j)} - \mu_j)^2}
\]

\subsubsection*{1. Step-by-Step Feature Mean and Standard Deviation Calculations}

\noindent\textbf{Feature 1: Living Area ($x^{(1)}$)}
\begin{align*}
    \mu_1 &= \frac{1.5 + 2.2 + 2.0 + 1.8 + 2.5 + 1.2 + 2.1 + 3.0 + 1.9 + 2.0 + 2.3 + 1.7}{12} \\
          &= \frac{24.2}{12} \\
          &\approx 2.0167 \\[6pt]
    \hspace*{1em}\sum_{i=1}^{12} (x_i^{(1)} - \mu_1)^2 & \\[4pt]
                &= (-0.5167)^2 + (0.1833)^2 + (-0.0167)^2 + (-0.2167)^2 + (0.4833)^2 + (-0.8167)^2 \\
                &\quad + (0.0833)^2 + (0.9833)^2 + (-0.1167)^2 + (-0.0167)^2 + (0.2833)^2 + (-0.3167)^2 \\
                &= 0.2670 + 0.0336 + 0.0003 + 0.0470 + 0.2336 + 0.6670 \\
                &\quad + 0.0069 + 0.9669 + 0.0136 + 0.0003 + 0.0803 + 0.1003 \\
                &= 2.4167 \\[6pt]
    \sigma_1 &= \sqrt{\frac{2.4167}{12}} \\
             &= \sqrt{0.2014} \\
             &\approx 0.4488
\end{align*}

\noindent\textbf{Feature 2: Distance to City Center ($x^{(2)}$)}
\begin{align*}
    \mu_2 &= \frac{1.5 + 0.5 + 1.0 + 1.2 + 2.0 + 3.0 + 0.8 + 5.0 + 1.5 + 0.5 + 1.1 + 1.0}{12} \\
          &= \frac{19.1}{12} \\
          &\approx 1.5917 \\[6pt]
    \hspace*{1em}\sum_{i=1}^{12} (x_i^{(2)} - \mu_2)^2 &\\[4pt]
          &= (-0.0917)^2 + (-1.0917)^2 + (-0.5917)^2 + (-0.3917)^2 + (0.4083)^2 + (1.4083)^2 \\
          &\quad + (-0.7917)^2 + (3.4083)^2 + (-0.0917)^2 + (-1.0917)^2 + (-0.4917)^2 + (-0.5917)^2 \\
          &= 0.0084 + 1.1917 + 0.3501 + 0.1534 + 0.1667 + 1.9834 \\
          &\quad + 0.6267 + 11.6168 + 0.0084 + 1.1917 + 0.2417 + 0.3501 \\
          &= 17.8892 \\[6pt]
    \sigma_2 &= \sqrt{\frac{17.8892}{12}} \\
             &= \sqrt{1.4908} \\
             &\approx 1.2210
\end{align*}

\noindent\textbf{Feature 3: Age of Property ($x^{(3)}$)}
\begin{align*}
    \mu_3 &= \frac{12 + 8 + 9 + 10 + 5 + 20 + 11 + 2 + 10 + 15 + 10 + 8}{12} \\
          &= \frac{120}{12}\\
          &= 10 \\[6pt]
    \hspace*{1em}\sum_{i=1}^{12} (x_i^{(3)} - \mu_3)^2 & \\[4pt]
            &= (2)^2 + (-2)^2 + (-1)^2 + (0)^2 + (-5)^2 + (10)^2 \\
            &\quad + (1)^2 + (-8)^2 + (0)^2 + (5)^2 + (0)^2 + (-2)^2 \\
            &= 4 + 4 + 1 + 0 + 25 + 100 \\
            &\quad + 1 + 64 + 0 + 25 + 0 + 4 = 228 \\[6pt]
    \sigma_3 &= \sqrt{\frac{228}{12}}\\
            &= \sqrt{19} \\
            &\approx 4.3589
\end{align*}

\noindent\textbf{Feature 4: Number of Bedrooms ($x^{(4)}$)}
\begin{align*}
    \mu_4 &= \frac{3 + 4 + 3 + 3 + 4 + 2 + 3 + 5 + 2 + 3 + 4 + 3}{12}\\
          &= \frac{39}{12} \\
          &= 3.2500 \\[6pt]
    \hspace*{1em}\sum_{i=1}^{12} (x_i^{(4)} - \mu_4)^2 & \\[4pt]
        &=(-0.25)^2 + (0.75)^2 + (-0.25)^2 + (-0.25)^2 + (0.75)^2 + (-1.25)^2 \\
        &\quad + (-0.25)^2 + (1.75)^2 + (-1.25)^2 + (-0.25)^2 + (0.75)^2 + (-0.25)^2 \\
        &= 0.0625 + 0.5625 + 0.0625 + 0.0625 + 0.5625 + 1.5625 \\
        &\quad + 0.0625 + 3.0625 + 1.5625 + 0.0625 + 0.5625 + 0.0625 \\
        &= 8.2500 \\[6pt]
    \sigma_4 &= \sqrt{\frac{8.2500}{12}}\\
             &= \sqrt{0.6875}\\
             & \approx 0.8292
\end{align*}

\subsubsection*{2. Normalizing the Query Point $\MVec{x}_q$}
Using our calculated parameters ($\mu_j, \sigma_j$) with $N=12$, we standardize the query vector $\MVec{x}_q = (2.0, 1.0, 10, 3)$:
\[
    x_{q,\text{scaled}}^{(1)} = \frac{2.0 - 2.0167}{0.4488} \approx -0.0372
\]
\[
    x_{q,\text{scaled}}^{(2)} = \frac{1.0 - 1.5917}{1.2210} \approx -0.4846
\]
\[
    x_{q,\text{scaled}}^{(3)} = \frac{10 - 10}{4.3589} = 0
\]
\[
    x_{q,\text{scaled}}^{(4)} = \frac{3 - 3.2500}{0.8292} \approx -0.3015
\]
\[
    \MVec{x}_{q,\text{scaled}} = (-0.0372, -0.4846, 0, -0.3015)
\]

\subsubsection*{3. Standardized Feature Matrix \& Updated Distances}
Applying $x_{\text{scaled}} = \frac{x - \mu}{\sigma}$ to the dataset produces the standardized feature matrix and recalculated Euclidean distances to $\MVec{x}_{q,\text{scaled}} = (-0.0372, -0.4846, 0, -0.3015)$:

\[
    \begin{array}{c|cccc|c|c}
        \text{Property} &
            x_{\text{scaled}}^{(1)} &
            x_{\text{scaled}}^{(2)} &
            x_{\text{scaled}}^{(3)} &
            x_{\text{scaled}}^{(4)} &
            y_i \text{ (Price)} &
            \mathbf{dist}(\MVec{x}_{q,\text{scaled}}, \MVec{x}_{i,\text{scaled}}) \\ \hline
        \MVec{x}_1 & -1.1512 & -0.0751 & 0.4588  & -0.3015 & 3.5 & 1.2725 \\
        \MVec{x}_2 & 0.4084  & -0.8941 & -0.4588 & 0.9045  & 5.2 & 1.4252 \\
        \MVec{x}_3 & -0.0372 & -0.4846 & -0.2294 & -0.3015 & 4.8 & \mathbf{0.2294} \quad (\text{Rank } 1) \\
        \MVec{x}_4 & -0.4828 & -0.3208 & 0       & -0.3015 & 4.6 & \mathbf{0.4748} \quad (\text{Rank } 3) \\
        \MVec{x}_5 & 1.0768  & 0.3344  & -1.1471 & 0.9045  & 6.0 & 2.1638 \\
        \MVec{x}_6 & -1.8197 & 1.1534  & 2.2942  & -1.5075 & 2.5 & 3.5466 \\
        \MVec{x}_7 & 0.1857  & -0.6484 & 0.2294  & -0.3015 & 4.9 & \mathbf{0.3594} \quad (\text{Rank } 2) \\
        \MVec{x}_8 & 2.1910  & 2.7912  & -1.8353 & 2.1105  & 7.0 & 4.9882 \\
        \MVec{x}_9 & -0.2600 & -0.0751 & 0       & -1.5075 & 4.1 & 1.2930 \\
        \MVec{x}_{10}& -0.0372 & -0.8941 & 1.1471  & -0.3015 & 4.3 & 1.2180 \\
        \MVec{x}_{11}& 0.6313  & -0.4027 & 0       & 0.9045  & 5.4 & 1.3813 \\
        \MVec{x}_{12}& -0.7056 & -0.4846 & -0.4588 & -0.3015 & 4.4 & 0.8107
    \end{array}
\]

\subsubsection*{4. Re-evaluating $K=3$ Neighbors and Final Prediction}
The 3 nearest neighbors under the standard deviation ($N=12$) normalization are:
\[
    \begin{array}{c|c|c}
        \text{Neighbor} & \text{Target } y_i & \mathbf{dist}(\MVec{x}_{q,\text{scaled}}, \MVec{x}_{i,\text{scaled}}) \\ \hline
        \MVec{x}_3 & 4.8 & 0.2294 \\
        \MVec{x}_7 & 4.9 & 0.3594 \\
        \MVec{x}_4 & 4.6 & 0.4748
    \end{array}
\]

The standardized KNN regression prediction yields:
\[
    \hat{y}(\MVec{x}_q) = \frac{4.8 + 4.9 + 4.6}{3} = \frac{14.3}{3} \approx 4.77 \text{ Crores}
\]

\subsection*{Key Insights}

\begin{tcolorbox}[colback=blue!5!white,colframe=blue!75!black,title=Classroom Takeaways: Why Feature Scaling Matters]
\begin{enumerate}
    \item \textbf{Equalizing Feature Weight (Standardization):} Using standard deviation ($\sigma = \sqrt{\frac{1}{N}\sum(x_i - \mu)^2}$) to rescale each feature column ensures that, after transformation, that feature has standard deviation $1$ across the dataset. This prevents higher-magnitude attributes (such as property age in years) from dominating the distance metric purely due to measurement units.
    \item \textbf{Shift in Neighborhood Rankings:}
          \begin{itemize}
              \item \textbf{Unstandardized Top Neighbor:} $\MVec{x}_4$ ($\text{dist} = 0.2828$), selected primarily because its property age matched $\MVec{x}_q$ perfectly ($10$ years), masking significant differences in living area and distance.
              \item \textbf{Standardized Top Neighbor:} $\MVec{x}_3$ ($\text{dist} = 0.2294$), which moves to Rank \#1 because all $4$ features contribute proportionally once mean-centered and scaled by $\sigma$.
          \end{itemize}
    \item \textbf{Generalization Rule:} Test/query points must always be normalized using the \textbf{training set} parameters ($\mu_{\text{train}}, \sigma_{\text{train}}$), never parameters recomputed from the test data — doing so causes information leakage.
\end{enumerate}
\end{tcolorbox}

\end{document}